%%%%%%%%%%%%%%%%%%%%%%%%%%%%%%%%%%%%%%%%%%%%%%%%%%%%%%%%%%%%%%%
%
% Welcome to Overleaf --- just edit your LaTeX on the left,
% and we'll compile it for you on the right. If you open the
% 'Share' menu, you can invite other users to edit at the same
% time. See www.overleaf.com/learn for more info. Enjoy!
%
%%%%%%%%%%%%%%%%%%%%%%%%%%%%%%%%%%%%%%%%%%%%%%%%%%%%%%%%%%%%%%%
%%%%%%%%%%%%%%%%%%%%%%%%%%%%%%%%%%%%%%%%%%%%%%%%%%%%%%%%%%%%%%%%%%%%%%%%
%% Skip to the "\begin{document}" line below to see the visible text. %%
%% Compare the left and right-hand panes to compare what you type and %%
%% how it appears when typeset.                                       %%
%%                                                                    %%
%% ***** Type answers after the next line that looks like this  ***** %%
%%%%%%%%%%%%%%%%%%%%%%%%%%%%%%%%%%%%%%%%%%%%%%%%%%%%%%%%%%%%%%%%%%%%%%%%
\documentclass[leqno,12pt]{article}
\usepackage{amsmath,amssymb}
\usepackage{graphicx}

\newcommand{\Term}{Spring 2026}
\newcommand{\Course}{CHM1046}

\newcommand{\Assignment}{Notes For Unit 01 Test}
\newcommand{\DueDate}{Friday, September~6}

\usepackage[body={6in,9in}]{geometry}

\newcommand{\brak}[1]{\left\langle #1\right\rangle}
\newcommand{\abs}[1]{\left\lvert #1\right\rvert}

\DeclareMathOperator{\lcm}{lcm}

\newcommand{\Neg}{\phantom{-}}
\renewcommand{\emptyset}{\varnothing}
\renewcommand{\subset}{\subseteq}

\newcommand{\ds}{\displaystyle}

%% Letter symbols
\newcommand{\Number}[1]{\mathbf{#1}}
\newcommand{\C}{\Number{C}}
\newcommand{\N}{\Number{N}}
\newcommand{\Q}{\Number{Q}}
\newcommand{\R}{\Number{R}}
\newcommand{\Z}{\Number{Z}}

\newcommand{\eps}{\varepsilon}

\newcommand{\Input}[1]{\includegraphics{#1.pdf}}

\newcommand{\Prob}[1]{\noindent\textbf{Question #1:}\quad}
\newcommand{\soln}[1][Answer]{\noindent\textbf{#1}\quad}

\newcommand{\Cmd}[1]{\textbackslash\texttt{#1}}

\pagestyle{empty}

\begin{document}
Aspen J. Johnson
\begin{center}

\textbf{Palm Beach State College - \Term} \\
\textbf{\Course} (Prof Marina Rines) \\
\textbf{\Assignment, Due \DueDate}
\end{center}

For \Course\ I'll be using the mathematical typesetting software
\LaTeX to code the following formula sheet, which will not be utilized during the exam. 


%%%%%%%%%%%%%%%%%%%%%%%%%%%%%%%%%%%%%%%%%%%%%%%%%%%%%%%%%%%%%%%%%%%%%%%%
%% *****      Type your answers after the "\soln" commands      ***** %%
%%%%%%%%%%%%%%%%%%%%%%%%%%%%%%%%%%%%%%%%%%%%%%%%%%%%%%%%%%%%%%%%%%%%%%%%
\begin{enumerate}

\item This item is automatically numbered. What does that tell you
  about \Cmd{item}?
\soln It tells me that \Cmd{item} will automatically number the item in numerical order. 

\item The double angle formula for the cosine function reads
  $\cos(2\theta) = \cos^{2}\theta - \sin^{2}\theta$.

  (a) Typeset the double angle formula for sine. (Look up this
  identity if you've forgotten it.)

  \soln $\sin(2\theta) = 2sin(\theta)\cos(\theta)$

  (b) Note that \Cmd{theta} gives the Greek letter~$\theta$. Typeset
  at least six other letters from the lowercase Greek alphabet. (Be
  sure your macros are in math mode!)

  \soln $\phi\delta\gamma\psi\beta\alpha$

  (c) Do you have a guess about how to get capital Greek letters?
  (Hint: \Cmd{Alpha} would be logical; as it happens, however, a
  capital alpha is a Roman capital~$A$, so there is no separate
  macro.) Typeset at least four uppercase Greek letters that have no
  Roman equivalent.

  \soln $\Gamma\Delta\Sigma\Phi$

\item If $y = \phi(x)$, and if we write $y + \Delta y = \phi(x + \Delta x)$,
  then
  \[ 
  \frac{\Delta y}{\Delta x} = \frac{\phi(x + \Delta x) - \phi(x)}{\Delta x},
  \qquad
  \frac{dy}{dx} = \lim_{\Delta x \to 0} \frac{\Delta y}{\Delta x}.
  \] 
  Modify parts of this code and see what happens. Things to watch for:
  What happens if you forget the backslash on a macro name? What
  happens if you typeset the line above using dollar signs (inline
  math mode) instead of display math? What happens if you add or
  remove space characters inside math mode?

\soln If you do not include the backlash on a macro name, the Greek letter is spelled out in English instead of being written as a Greek symbol. If you use dollar signs instead of display math, it deletes the previously existing break in the text allotted to the equation. It appears as though nothing changes when you add or remove characters inside math mode.  
\item If $a$, $b$, and~$c$ are non-zero numbers, and if
  $ax^{2} + bx + c = 0$, typeset the (quadratic) formula for~$x$ in
  terms of $a$,~$b$, and~$c$.

  Hint: The macro \Cmd{pm} gives a plus-or-minus sign.\footnote{The
    tilde character is a \emph{tie} or \emph{non-breaking
      space}. Using ties is a little fussy, but helps distinguish
    professional typography. Use a tie where a line break would
    interrupt the reader's attention.}

  \soln \[
  x=\frac{-b\pm\sqrt{b^2-4ac}}{2a}
\]
\item One fundamental theorem of calculus asserts that if $f$~is a
  continuous function on a closed, bounded interval $[a, b]$, then
  \[
  \frac{d}{dx} \int_{a}^{x} f(t)\, dt \qquad = f(x)
  \quad\text{for all~$x$ with $a \leq x \leq b$.}
  \]
  (a) Modify parts of this code and see what happens. Hints:
  \Cmd{,}~gives a thin space, \Cmd{quad} gives one em of space, the
  width of the letter~m in the current font, \Cmd{qquad} gives two
  ems. The underscore and caret each accept one argument, and make a
  subscript or superscript. The \Cmd{text} command sets its argument
  in ``paragraph mode''. Why is it needed here?
  
  \soln Without \Cmd{text}, the words "for all x with..." all become once word and the computer reads it as part of the math equation instead of as English. It also allows for the break in the text. 
  (b) Evaluate the following derivatives and typeset the answers:
  \[
  \frac{d}{dx} \int_{0}^{x} e^{\sin t}\, dt;\qquad
  \frac{d}{du} \int_{\pi}^{u} \log(1 + e^{t^{2}})\, dt;\qquad
  \frac{d}{dt} \int_{-1}^{t} \frac{\sqrt{\tan s}}{1 + s^{2}}\, ds.
  \]

  \soln \[
  e^{\sin x} \qquad \log(1+e^{u^{2}})\ \qquad \frac{\sqrt{tan t}}{1+t^{2}}\ 
\]
  (c) Typeset the other fundamental theorem of calculus, sometimes
  known as the \emph{evaluation theorem}. Don't forget to state
  appropriate hypotheses.

  \soln If \emph{f} is continuous at every point of [\emph{a,b}], and if \emph{F} is any anti derivative of \emph{f} on [\emph{a,b}], then 
  \[
\int_{a}^{b} f(x)\, dx = F(b)-F(a)
\]
\end{enumerate}

\end{document}